\section{Hoja de ruta}
\label{sec:fases}
\begin{respuesta}
Empezaría por la visibilidad: en dos semanas, un gateway frente a los agentes existentes responde
cuántos hay, con qué modelo, qué datos y cuánto cuestan. Luego el núcleo con un agente real
(facturación). El autoservicio llega cuando la compuerta existe. Multiagente, enrutamiento
dinámico y fine-tuning, solo con señal. La regla: \emph{nunca abrir una puerta antes de tener
el control que la vigila}.
\end{respuesta}

\begin{figure}[!htbp]
  \centering
  \begin{ganttchart}[
    hgrid, vgrid={*{3}{draw=none},dotted},
    x unit=0.53cm, y unit chart=0.5cm, y unit title=0.5cm,
    title height=1, title label font=\sffamily\tiny,
    bar label font=\sffamily\scriptsize, group label font=\sffamily\scriptsize\bfseries,
    milestone label font=\sffamily\scriptsize\itshape\color{docacento},
    bar/.append style={fill=execbg, draw=execbd},
    group/.append style={fill=docink, draw=docink},
    milestone/.append style={fill=docacento, draw=docacento},
    bar height=0.55, group peaks height=0.2, group peaks tip position=0,
    bar label node/.append style={align=right},
  ]{1}{26}
  \gantttitle{Semanas}{26}\\
  \gantttitlelist{1,...,26}{1}\\
  \ganttgroup{Fase 0 · Visibilidad}{1}{2}\\
  \ganttbar{Inventario + gateway}{1}{2}\\
  \ganttmilestone{Vista de conjunto}{2}\\
  \ganttgroup{Fase 1 · Núcleo}{3}{8}\\
  \ganttbar{Manifiesto, registro, runtime}{3}{6}\\
  \ganttbar{MCP facturación y medición}{4}{7}\\
  \ganttbar{Suites y compuerta v1}{5}{8}\\
  \ganttbar{Migración facturación}{6}{8}\\
  \ganttmilestone{Primer agente certificado}{8}\\
  \ganttgroup{Fase 2 · Autoservicio}{9}{16}\\
  \ganttbar{Constructor + Arquitecto}{9}{14}\\
  \ganttbar{Conocimiento con vigencia}{9}{12}\\
  \ganttbar{Evaluación en línea}{11}{16}\\
  \ganttmilestone{Autoservicio abierto}{16}\\
  \ganttgroup{Fase 3 · Escala}{17}{26}\\
  \ganttbar{Flujo de cobranzas (L4)}{17}{22}\\
  \ganttbar{Enrutamiento por costo}{19}{24}\\
  \ganttbar{Fine-tuning (si hay señal)}{22}{26}
\end{ganttchart}

  \caption{Cronograma grueso en semanas. Los rombos son los hitos que cambian lo que la
  organización puede hacer.}
  \label{fig:gantt}
\end{figure}

\begin{table}[!htbp]
\centering\small
\caption{Fases: qué entra, qué se adopta y el riesgo que más preocupa.}
\label{tab:fases}
\renewcommand{\arraystretch}{1.12}
\begin{tabularx}{\textwidth}{@{}L{2.4cm}YL{2.6cm}L{3.8cm}@{}}
\toprule
\textbf{Fase} & \textbf{Qué entra} & \textbf{Se adopta} & \textbf{Riesgo y mitigación}\\
\midrule
0 · Visibilidad\newline sem.\ 1--2 & Inventario de agentes actuales; gateway con llaves centralizadas; costo por área & LiteLLM, Langfuse & Resistencia de áreas: solo cambian URL y llave\\
1 · Núcleo\newline sem.\ 3--8 & Manifiesto, registro, runtime con guards, MCP de facturación y medición, compuerta v1, migrar facturación & LangGraph, pgvector & Sobrediseñar: se diseña contra un agente real\\
2 · Autoservicio\newline sem.\ 9--16 & Constructor y Arquitecto, ingesta con vigencia, evaluación en línea, caducidad & SSO existente & Agentes mediocres: compuerta, dueño y caducidad\\
3 · Escala\newline sem.\ 17--26 & Cobranzas L4, enrutamiento por costo, cobro por área, fine-tuning si hay señal & --- & Complejidad multiagente: traspasos tipados\\
\bottomrule
\end{tabularx}
\end{table}

\decision{Construir solo lo específico del negocio; adoptar el resto.}
{Se construyen manifiesto, registro, constructor, compuerta y conectores MCP: es el
conocimiento que ningún proveedor tiene. Se adoptan gateway, runtime, vectores,
observabilidad y SSO.}
{Comprar una plataforma de agentes completa: más rápida en la fase 2, pero ata modelos y deja
los inventarios fuera de la empresa.}

\textbf{El riesgo que más me preocupa} es de confianza: que un agente exponga datos de un cliente
a otro o ejecute una acción no autorizada. Por eso sus tres defensas (sujeto desde la sesión,
escritura con confirmación, guard de salida por bloques) están en la fase 1, antes del primer
agente publicado. Es coherente con el NIST AI RMF \citep{nist2023rmf}, la ISO/IEC 42001
\citep{iso42001} y el eje de mitigación de riesgos del CONPES 4144 \citep{conpes4144}.
